\documentclass[10pt,a4paper]{amsart}
\usepackage[margin=19mm]{geometry}
\usepackage{amsmath,amssymb,mathtools}
\usepackage[scr=rsfs,cal=euler]{mathalfa}
\usepackage{xfrac}
\usepackage[unicode,hidelinks,bookmarks=false]{hyperref}
\newcommand{\ie}{i.e.\ }
\newcommand{\cf}{cf.\ }
\newcommand{\resp}{respectively\ }
\newcommand{\Subsection}{\subsection}
\providecommand{\nobreakdash}{\nobreak\hbox{-}}
\setlength{\emergencystretch}{2em}
\multlinegap0pt
\usepackage{bm}
\usepackage{bbm}
\usepackage{dsfont}
\usepackage{xspace}
\usepackage{cancel}
\usepackage{mathtools}
\usepackage{amsthm}
\makeatletter
\@ifundefined{theorem}{\newtheorem{theorem}{Theorem}}{}
\makeatother
\title[On Intersection and Co-maximal Hypergraph of $\mathbb{Z}\_n$]{On Intersection and Co-maximal Hypergraph of $\mathbb{Z}\_n$}
\author{Sachin Ballal, Ardra A N}
\date{Source: arXiv:2509.14618v1 (CC BY 4.0)}
\begin{document}
\maketitle

\noindent\emph{Source and licence.} On Intersection and Co-maximal Hypergraph of $\mathbb{Z}_n$. Version arXiv:2509.14618v1 (CC BY 4.0), distributed under the Creative Commons Attribution 4.0 International licence, as recorded in the arXiv licence metadata for this exact version and shown on the abstract page https://arxiv.org/abs/2509.14618v1. Full rights notice: licenses/arxiv-2509.14618-CC-BY-4.0.txt.\par\smallskip

\noindent\textit{Editorial scope.} Counted unit: the Theorem whose source label is \texttt{iso} in the article (source0.tex lines 128--130, source label \texttt{iso}), delivered together with the article's own complete original proof of it (source0.tex lines 131--167, one \texttt{proof} environment of 37 lines). No mathematics is written, repaired or re-derived: statement and proof are the article's own text line for line. Required context retained uncounted: none. Every definition, display and label of those lines is retained unchanged. Omitted from this package and not redistributed: 1--127, 168--2416. Nothing inside a retained excerpt is omitted or altered: every retained body is delivered exactly as the article's lines are written, so no omission receipt of any kind accompanies this package. Citations: the retained excerpts carry no citation command at all, so no bibliography entry of the article is transcribed and no reference is redistributed. Reference adaptation: none was needed and none was made; every reference inside the delivered unit resolves inside it, so no unresolved reference or citation is produced. Printed slips kept literal: no printed word, symbol, hypothesis or number of the counted statement or of its proof was altered. Layout and class: the article's own class is \texttt{article} with packages including geometry, fontenc, tikz, inputenc, graphicx, mathtools, amssymb, amsthm, thmtools, xcolor, nameref, hyperref, amsmath, amsfonts; the wrapper is the lane's pinned \texttt{amsart} editorial wrapper at 10pt on A4, which restates the theorem-like environments the retained text needs and adds the article's own macro definitions that the retained text needs (none), with extra wrapper packages (bm, bbm, dsfont, xspace, cancel, mathtools). The delivered numbering is the unit's own. Assets: no retained span contains a figure environment, a graphics-inclusion command, a PDF-page inclusion, a table or a TikZ picture; no image file is copied into this package, the package declares no asset dependency and no image file is redistributed with it. Licence: version arXiv:2509.14618v1 of this article is distributed under the Creative Commons Attribution 4.0 International licence, as recorded in the arXiv licence metadata for this exact version and shown on the abstract page https://arxiv.org/abs/2509.14618v1. A full rights notice is delivered at licenses/arxiv-2509.14618-CC-BY-4.0.txt. Characters: every retained excerpt is pure ASCII, so the delivered engine, XeLaTeX with its default Latin Modern text font, reports no missing character.
	\begin{theorem}\label{iso}
		For a finite abelian group $G, \tilde{\Gamma}_{\mathcal{H}}(G) \cong Co_\mathcal{H}(G)$. 
	\end{theorem}

	\begin{proof}
		Let $G$ be a finite abelian group and $H_1 \in V(\tilde{\Gamma}_{\mathcal{H}}(G))$.
		Then, by the definition, there exists $H_2  \in V(\tilde{\Gamma}_{\mathcal{H}}(G))$ such that $H_1 \cap H_2 =\{e\}$. Note that for subgroups $H_1,H_2$ of $G$, $H_1 \vee H_2=$ $<H_1 \cup H_2>$ and $H_1 \wedge H_2 = H_1 \cap H_2$ in $L(G)$, the lattice of subgroups of $G$. As $G$ is abelian, $\delta(H_1)$ and $\delta(H_2)$ are normal. Therefore, 
		$\delta(H_1) \delta(H_2) = \delta(H_1) \vee \delta(H_2) =
			  \rho \perp(H_1) \vee \rho\perp(H_2) = \\ \rho(\perp(H_1) \vee \perp(H_2))   = \rho(\perp(H_1 \wedge H_2)) 
			  = \rho((H_1 \wedge H_2)^{\perp})
			  = \rho(\{e\}^{\perp})
			 = \rho (G^\star)
			 = G$
		Thus, if $H_1 \in V(\tilde{\Gamma}_{\mathcal{H}}(G))$, then $\delta(H_1) \in V(Co_\mathcal{H}(G))$.\\
		Define a map $f:    V(\tilde{\Gamma}_{\mathcal{H}}(G)) \longrightarrow V(Co_\mathcal{H}(G))$ such that $f(H) = \delta(H)$ for each $H \in  V(\tilde{\Gamma}_{\mathcal{H}}(G))$. Since $\delta$ is a well-defined bijective map, $f$ is also a well-defined bijective map.\\ Observe that,  
		$
		H_1 \sim H_2 \,\, \text{in} \, \, \tilde{\Gamma}_{\mathcal{H}(G)} \, \, 
			 \text{if and only if} \, \,  H_1 \cap H_2 =\{e\}
			 \text{if and only if} \, \,(H_1 \cap H_2)^{\perp} = \{e\}^{\perp}
		\, \, 
		 \text{if and}\\ \text{only if} \, \, (H_1 \wedge H_2)^{\perp} = G^\star 
		\, \, 
		 \text{if and only if} \, \, \rho((H_1 \wedge H_2)^{\perp}) = \rho(G^\star)
		\, \, 
		 \text{if and only if} \, \, \rho(\perp(H_1 \wedge H_2)) = G 
			\, \, 
			 \text{if and only if} \, \,\rho(\perp(H_1) \vee \perp(H_2)) = G
			\, \, 
			 \text{if and only if} \, \, \rho \perp(H_1) \vee \rho\perp(H_2) = G 
		\, \, 
		 \text{if and only if}\\ \, \, \delta(H_1) \vee \delta(H_2) = G 
		\, \, 
		 \text{if and only if} \, \,\delta(H_1)  \delta(H_2) = G 
			\, \, 
			 \text{if and only if} \, \, f(H_1)  f(H_2) = G 
		\, \, 
		 \text{if and only if}\\ \, \, f(H_1) \sim f(H_2) \,\, \text{in} \,\, Co_\mathcal{H}(G)$, i.e., $
		 H_1 \sim H_2 \,\, \text{in} \, \, \tilde{\Gamma}_{\mathcal{H}(G)}$ if and only if $f(H_1) \sim f(H_2) \,\, \text{in} \,\, Co_\mathcal{H}(G)$. 
		Hence,  for a positive integer $r$,  $\{H_1, H_2, \ldots, H_r\}$ is a hyperedge  of $\tilde{\Gamma}_{\mathcal{H}}(G)$ if and only if \\ $\{f(H_1),f(H_2), \dots, f(H_r)\}$ is a hyperedge of $Co_\mathcal{H}(G)$.\\
		Thus, $f$ is an isomorphism from $V(\tilde{\Gamma}_{\mathcal{H}}(G))$ to $V(Co_\mathcal{H}(G))$. Thus, $\tilde{\Gamma}_{\mathcal{H}}(G) \cong Co_\mathcal{H}(G)$. 
	\end{proof}

\end{document}
